\section*{§1. The Field of Complex Numbers}

The set of complex numbers is denoted by $\mathbb{C}$. A complex number is an ordered pair $(a,b)$ of real numbers. We define addition and multiplication of complex numbers by the formulas
\begin{align*}
(a,b)+(c,d)&=(a+c,b+d)\\
(a,b)\cdot(c,d)&=(ac-bd,ad+bc).
\end{align*}
It is easy to check that these operations satisfy the axioms of a field. The zero element is $(0,0)$, the unit element is $(1,0)$, and the inverse of $(a,b)$ is
\[
(a,b)^{-1}=\left(\frac{a}{a^{2}+b^{2}},\frac{-b}{a^{2}+b^{2}}\right).
\]
The field of complex numbers contains the field of real numbers as a subfield, since the set of pairs of the form $(a,0)$ is a subfield of $\mathbb{C}$ isomorphic to $\mathbb{R}$. We shall identify the real number $a$ with the complex number $(a,0)$. In particular, $0=(0,0)$ and $1=(1,0)$. The complex number $(0,1)$ is denoted by $i$. Then,
\[
i^{2}=(0,1)\cdot(0,1)=(-1,0)=-1,
\]
so that $i$ is a square root of $-1$. Every complex number $(a,b)$ can be written in the form
\[
(a,b)=(a,0)+(0,b)=(a,0)+(0,1)\cdot(b,0)=a+ib.
\]
Thus, every complex number can be written in the form $a+ib$, where $a$ and $b$ are real numbers. The number $a$ is called the \textbf{real part} of the complex number $a+ib$, and $b$ is called the \textbf{imaginary part}. We write
\[
a=\text{Re}(a+ib),\quad b=\text{Im}(a+ib).
\]
The complex number $a-ib$ is called the \textbf{conjugate} of the complex number $a+ib$, and is denoted by $\overline{a+ib}$. We have
\[
\overline{z+w}=\overline{z}+\overline{w},\quad \overline{zw}=\overline{z}\cdot\overline{w},\quad \overline{\overline{z}}=z
\]
for any complex numbers $z$ and $w$. The \textbf{absolute value} or \textbf{modulus} of a complex number $z=a+ib$ is defined by
\[
|z|=\sqrt{a^{2}+b^{2}}.
\]
We have
\[
z\overline{z}=|z|^{2},\quad |zw|=|z|\cdot|w|
\]
for any complex numbers $z$ and $w$.